\documentclass[11pt,letterpaper]{article}
\usepackage[T1]{fontenc}
\usepackage[utf8]{inputenc}
\usepackage{lmodern,microtype}
\usepackage[margin=0.88in,headheight=15pt]{geometry}
\usepackage{amsmath,amssymb,amsthm,mathtools}
\usepackage{booktabs,longtable,array,tabularx}
\usepackage{xcolor,listings,fancyhdr,xurl,enumitem}
\usepackage{hyperref,bookmark}
\definecolor{ink}{HTML}{17354A}
\definecolor{accent}{HTML}{176B72}
\definecolor{paper}{HTML}{F0F5F6}
\hypersetup{colorlinks=true,linkcolor=ink,urlcolor=accent,citecolor=ink,
 pdftitle={AsymptoticAnalysis: validating core components before cancellation},
 pdfauthor={OpenAI},pdfsubject={Incremental source audit at efa1aeec4845a9c35e140963a0333d0c9ec33b05}}
\pagestyle{fancy}\fancyhf{}
\fancyhead[L]{\small\sffamily ASYMPTOTICANALYSIS / COMPONENT VALIDATION}
\fancyhead[R]{\small\sffamily 10 September 2026}
\fancyfoot[L]{\small\sffamily Incremental audit \textbullet\ Wolfram 15 / Mathics3}
\fancyfoot[R]{\small\sffamily\thepage}
\setlength{\parindent}{0pt}
\setlength{\parskip}{5.5pt plus 1pt}
\setlength{\emergencystretch}{2em}
\setlist{nosep,leftmargin=1.6em}
\lstset{basicstyle=\ttfamily\footnotesize,breaklines=true,columns=fullflexible,
 keepspaces=true,showstringspaces=false,frame=single,rulecolor=\color{black!15},
 backgroundcolor=\color{paper},xleftmargin=4pt,xrightmargin=4pt,
 aboveskip=9pt,belowskip=9pt,tabsize=2}
\newtheorem{theorem}{Theorem}[section]
\newtheorem{proposition}[theorem]{Proposition}
\newtheorem{lemma}[theorem]{Lemma}
\theoremstyle{remark}\newtheorem{remark}[theorem]{Remark}
\newcommand{\code}[1]{\texttt{#1}}
\newcommand{\file}[1]{\path{#1}}
\newcommand{\OO}[1]{\mathcal O\!\left(#1\right)}
\newcommand{\RR}{\mathbb R}
\newcommand{\rev}{\texttt{efa1aee}}
\newcommand{\fullrev}{\texttt{efa1aeec4845a9c35e140963a0333d0c9ec33b05}}
\newcommand{\eps}{\varepsilon}
\newcolumntype{Y}{>{\raggedright\arraybackslash}X}

\begin{document}
\begin{titlepage}
\vspace*{0.4in}
{\sffamily\bfseries\color{accent} INCREMENTAL CODE AUDIT / AFTER SIX REVIEW WAVES\par}
\vspace{0.35in}
{\sffamily\bfseries\color{ink}\fontsize{31}{36}\selectfont AsymptoticAnalysis\par}
\vspace{0.23in}
{\sffamily\fontsize{23}{29}\selectfont Validating core components\newline before cancellation\par}
\vspace{0.35in}
{\large A reproduced remainder-contract failure, exact all-depth\newline counterexamples, and a focused repair for the shared source\par}
\vspace{0.45in}
\colorbox{paper}{\parbox{0.94\textwidth}{\vspace{10pt}
\textbf{Principal result.} The exact-core constructor checks whether
\code{core + perturbation} is target-free, then uses the two components
separately. Cancelling occurrences of \code{Re[y]} bypass that check. Two
successful public calls for the equation $x+x^{-1}=y$ claim errors of
$O(y^{-2})$ and $O(y^{-4})$, although both actual errors are $\Theta(y^{-1})$.
A component-wise guard rejects the offending split without forbidding a
legitimate target-dependent \code{CoreInverse} expression.\vspace{10pt}}}
\vfill
\begin{tabular}{@{}ll@{}}
Repository & \code{VladimirReshetnikov/Asymptotic}\\
Pinned revision & \fullrev\\
Revision time & 10 September 2026, 19:09:04 UTC\\
New finding families & One, with two public counterexamples\\
Native execution & Wolfram 15.0.0, Linux x86 (64-bit)\\
Mathics execution & Not performed; portable regression probe supplied\\
Independent checks & 30 Python tests, all passing
\end{tabular}\par
\vspace{0.25in}
{\small Scope: broad, incremental mechanism-level review with targeted execution.
This is not a full-suite acceptance record or a claim that every source line
and every archived report paragraph has been independently verified.}
\end{titlepage}

\begin{abstract}
This report audits selected analytic engines, coordinate transformations,
coefficient and remainder calculations, Mathics adapters, and result contracts
at a pinned revision of AsymptoticAnalysis. The exclusion set comprises the
repository's six review-wave indexes, maintained finding register and intakes,
and targeted original reports. After removing already recorded mechanisms and
ruling out several plausible but incorrect candidates, the report retains one
new high-priority correctness finding.

The defect is in the admission of the separate core and perturbation arguments
to \code{AsymptoticCoreInverse}. Validation of their evaluated sum loses target
occurrences that cancel, whereas the subsequent marker expansion retains those
occurrences in the two components. Consequently a fixed-data asymptotic
majorant is applied to moving data. Public Wolfram 15.0.0 executions reproduce
false remainder classes at both an infinite and a finite source endpoint.
Elementary quadratic inverses prove the errors exactly. A Catalan coefficient
formula extends the diagnosis to every fixed marker depth and explains why
marker convergence does not rescue the claimed power order.

A two-line candidate patch was exercised in a temporary, commit-pinned
standalone copy: five focused public requests exhibited the expected refusals
or preserved controls. Separately, 22 exact mathematical tests and eight
patch-utility fixture tests pass. A 17-case native forward-expansion screen
checks a different part of the package and is reported separately. The archive
contains the article, patch, exact rational oracles, numerical enclosures,
observation records, novelty ledger, source inventory, and unexecuted complete
Wolfram/portable regression files. No Mathics execution, Windows-kernel
acceptance, or full repository-suite execution is asserted.
\end{abstract}

\setcounter{tocdepth}{1}
\begingroup
\small\setlength{\parskip}{0pt}
\tableofcontents
\endgroup
\newpage

\section{Executive assessment}
\subsection{The retained new finding}
\textbf{N01 --- target dependence erased by sum validation.}
The constructor accepts a core decomposition that its own fixed-data remainder
theorem does not admit. The finite approximation and the remainder are both
part of the returned analytic object; checking only that the forward equation
simplifies to a target-free expression is insufficient. In the large-source
witness, the first inverse correction is wrong by $1/(4y)$ at marker depth two.
In the finite-source witness, even the leading coefficient is wrong: $7/8$
instead of $1$. Both objects nevertheless carry much smaller error classes.
The observed source mechanism is at lines 126--127 of
\file{src/Kernel/CorePerturbation.wl}; the downstream bound is formed in the
same module \cite{core}.

The recommended immediate change is to test the tuple
\code{\{core, perturbation\}}, not its evaluated sum. This is a local
correctness repair. It neither broadens the supported asymptotic scale nor
changes which parameters a user-supplied inverse is allowed to contain.
The tuple-based \code{FreeQ} check alone closes the reproduced target leak;
the accompanying tuple-based general input validation aligns the validation
object with the data actually consumed.

\subsection{Why one finding can justify a substantial incremental report}
The pinned tree indexes 55 numbered reports across six waves, with 45 retained
packages and ten retirement tombstones. Numerous earlier correctness,
performance, API, and portability issues are already recorded. Those are not
repackaged here as new discoveries \cite{reviewindex,register,wave3,wave4,wave5,wave6}.
The useful deliverable is therefore a sharply delimited new failure, with a
proof, reproducible evidence, a small repair, and acceptance obligations that
are specific to that repair.

The source examination was broader than the retained finding. It included
ordinary forward expansions and observables; Fourier and logarithmic
coefficient paths; flat and exact-core sectors; Gamma/Barnes inverses;
Dirichlet and parameterized special-function paths; coordinate reconstruction;
interval-certificate logic; and selected Mathics proof/algebra/Taylor adapters.
The machine-readable inventory records whether a file was read completely or
only in relevant ranges. No conclusion of global correctness follows merely
because no additional unreported defect was retained.

\subsection{Decision summary}
\begin{center}
\begin{tabularx}{\textwidth}{@{}p{0.20\textwidth}Y@{}}
\toprule
Decision & Basis and boundary\\
\midrule
Repair now & N01 has two successful public reproductions, exact counterexamples,
and a component-local admission fix.\\
Keep the valid API & A fixed real core parameter remains admissible; a supplied
\code{CoreInverse} must still be allowed to depend on the target.\\
Do not infer Mathics acceptance & The source repair is shared, but no fresh
Mathics package execution was available. Earlier realness checks may alter
public reachability there.\\
Do not claim a performance defect & Known dense-allocation and recurrence
issues were excluded. No new measured end-to-end slowdown is established here.\\
Do not widen the theorem silently & Moving core parameters need a different
joint asymptotic argument, not just a successful symbolic marker calculation.\\
\bottomrule
\end{tabularx}
\end{center}

\section{Baseline, source coverage, and non-overlap}
\subsection{The exact reviewed artifact}
The repository revision is \fullrev, recorded at 19:09:04 UTC on
10 September 2026. The canonical implementation resides in
\file{src/Kernel/}; the root \file{AsymptoticAnalysis.wl} is the generated,
self-contained entry. Native probes downloaded that entry at the full commit
identifier rather than following a moving \code{main} URL
\cite{repo,sourcemap,guide}.

A local repository checkout was not obtained in this environment. Source was
read through the GitHub connector, and the native execution service downloaded
the pinned standalone file independently. This distinction matters:
reading source excerpts, loading the entire generated package, and exercising
selected public calls are three different activities. The successful native
load returned \code{Null} and produced the expected ordinary inverse smoke
result. No repository file, branch, or remote review directory was changed.

\subsection{Coverage by mechanism}
\begin{longtable}{@{}p{0.26\textwidth}p{0.68\textwidth}@{}}
\toprule
Area & Examination and outcome\\
\midrule\endhead
Exact-core inversion & Model construction, automatic inverse, marker recurrence,
input validation, output fields, and fixed-data majorant examined. N01 retained.\\
Ordinary forward calculus & Selected standalone definitions plus 17 native
coefficient comparisons. A varying-coefficient observable control transports
its lost precision correctly.\\
Fourier coefficient engine & Source ranges through construction, coefficient,
and residual APIs; finite-endpoint power convention checked against the guide.
That suspected power bug was ruled out.\\
Flat and logarithmic sectors & Finite flat model/recurrence/majorant and selected
logarithmic model/unit-recursion paths examined. Existing reviewed obligations
were not reissued.\\
Gamma/Barnes and special functions & Coefficient recurrences, frontiers,
parameterized Taylor/Frobenius ingress, selected real-domain rules, and
Dirichlet formulas inspected at the scope in the inventory.\\
Coordinate transformations & Target-log construction and selected source-chart
reconstruction and observable contracts inspected.\\
Exact certificates & Selected interval arithmetic and source-domain proof paths
examined. This report does not claim a new certificate soundness finding.\\
Mathics compatibility & Assumption, algebra, core-function, special-function,
and Taylor adapters inspected. Runtime success is not inferred from source
inspection.\\
Prior review corpus & Six wave indexes, maintained status/intakes, and targeted
nearest original reports used for mechanism-level exclusion. Not a word-for-word
rereading of every retained article.\\
\bottomrule
\end{longtable}
The 32-entry \file{evidence/source-inventory.csv} supplies pinned URLs and
file-specific inspection boundaries. Some broad files were fetched in full but
read in selected relevant sections; the inventory deliberately does not label
that as a complete line-by-line audit.

\subsection{Nearest earlier findings}
The closest earlier result is report 46's target-dependent
\code{SourceShift} in \code{Asymptotic\allowbreak Exponential\allowbreak CoreInverse}, also carried by
report 50. It applies the same essential mathematical principle---parameters
called fixed must really be fixed---but concerns a different public API,
option, validation line, normalized exponential amplitude, and repair
\cite{r46,wave6}. N01 uses no \code{SourceShift} and no nondefault option.
Its target occurrence is hidden by cancellation across two positional
arguments. The earlier option guard does not repair that test.

Report 43 F03 concerns a supplied inverse/formula retaining the eliminated
source variable. N01 uses automatic core inversion and supplies no
source-bearing \code{CoreInverse}. The C15 composition issue concerns a varying
parameter captured while combining stored series; N01 occurs before a single
exact-core result is constructed. Option-identity issues likewise do not
explain the option-free counterexample \cite{r43,register,wave3,wave6}.

These distinctions are recorded in \file{evidence/novelty-ledger.csv}. They are
not assertions that the fixed-data principle itself is novel mathematics.
The new contribution is a previously unmatched admission mechanism and its
public failure, together with the corresponding local remedy. An unindexed
paragraph elsewhere in the historical archive could warrant additional
attribution; the claimed exclusion scope is the reviewed indexes, registers,
intakes, and targeted source reports, not every archived sentence.

\section{N01: a wrong remainder at an infinite source endpoint}
\subsection{A minimal public request}
Load the pinned package in a fresh native kernel, then evaluate:
\begin{lstlisting}
s = AsymptoticAnalysis`AsymptoticCoreInverse[
  x - Re[y], Re[y] + 1/x,
  {x, Infinity}, {y, 2}
];
\end{lstlisting}
The two arguments add to
\begin{equation}\label{eq:original}
 f(x)=x+\frac1x.
\end{equation}
The target symbol has disappeared from this sum. In the successful Wolfram
15.0.0 execution, the result was a \code{GeneralizedSeries}, with the following
actual projections:
\begin{lstlisting}
Normal[s]
(* y - 1/(y + Re[y])
     - (Re[y] + 1/(y + Re[y]))/(y + Re[y])^2 *)

s["Remainder"]
(* PowerLogRemainder[1/(y + Re[y]), 2, 0] *)

s["RemainderScaleExpression"]
(* 1/(y + Re[y])^2 *)

s["Function"]
(* x + 1/x *)
\end{lstlisting}
The retained domain is
\begin{lstlisting}
y + Re[y] > 0 &&
Element[1/(y + Re[y]), Reals] &&
0 < 1/(y + Re[y]) < 1/E
\end{lstlisting}
Every condition holds on the real ray $y>3$. Thus the witness is not obtained
by evaluating outside the result's declared domain. The domain is a branch and
coordinate condition; it does not restore the missing fixed-parameter premise.

\subsection{Exact comparison with the selected inverse}
For $y>2$, the inverse branch tending to $+\infty$ is
\begin{equation}\label{eq:roots}
 g_+(y)=\frac{y+\sqrt{y^2-4}}2,\qquad
 g_-(y)=\frac{y-\sqrt{y^2-4}}2=\frac1{g_+(y)}.
\end{equation}
Both identities follow directly from $x^2-yx+1=0$. The chosen large root is
positive and exceeds one; the forward derivative $1-x^{-2}$ is positive
there, so there is no local branch ambiguity.

On $y>3$, the returned approximation and scale simplify to
\begin{equation}\label{eq:largeobserved}
 A_2(y)=y-\frac{3}{4y}-\frac{1}{8y^3},\qquad
 R_2^+(y)=\frac1{4y^2}.
\end{equation}
The exact root expansion is
\begin{equation}\label{eq:catalanroot}
 g_+(y)=y-\frac1y-\frac1{y^3}-\frac2{y^5}-\frac5{y^7}-\cdots.
\end{equation}
Consequently,
\begin{equation}\label{eq:largeerror}
 A_2(y)-g_+(y)=\frac1{4y}+\frac7{8y^3}+\OO{y^{-5}}.
\end{equation}
In particular,
\begin{equation}\label{eq:largeratio}
 \frac{|A_2(y)-g_+(y)|}{R_2^+(y)}\sim y\longrightarrow+\infty.
\end{equation}
No fixed multiplicative constant makes the advertised remainder class true
on a sufficiently far positive ray.

This is stronger than a numerical discrepancy at a moderate target. It is an
asymptotic contradiction. Nor is it merely a conservative bound: a conservative
bound would be larger than necessary, whereas this one is too small by an
unbounded factor. The displayed finite expression is also not the correct
ordinary inverse expansion through the advertised precision.

\subsection{Why \code{Re[y]} is significant}
The first attempted request with \code{x-y} and \code{y+1/x} returned
\code{Failure["UnprovedCoreData",\ldots]}. That is a negative candidate, not a
public wrong-result witness. Replacing the cancelling target occurrences by
\code{Re[y]} exposes the actual missing condition: the realness test can admit
the coefficient even though it is not fixed on the target approach.
The function \code{Re} expresses real part; it does not assert independence
from the symbol inside it \cite{re}.

The point is not to ask \code{Assumptions} to declare the target a parameter.
The constructor already rejects source/target occurrences in its parameter
assumptions. Rather, this request makes target dependence survive in data
whose reality is otherwise acceptable. A test for real coefficients is not a
substitute for the target-free test.

\section{The same defect at a finite source endpoint}
\subsection{A second public manifestation, not a second finding}
The complementary source chart gives an even sharper witness:
\begin{lstlisting}
t = AsymptoticAnalysis`AsymptoticCoreInverse[
  1/x - Re[y], Re[y] + x,
  {x, 0}, {y, 2}
];
\end{lstlisting}
Again the forward function is $x+x^{-1}$. Here the source tends to $0^+$ and
the correct inverse is $g_-(y)$ from \eqref{eq:roots}. The successful native
result, simplified under $y>3$, is
\begin{equation}\label{eq:smallobserved}
 B_2(y)=\frac{1+5y^2+14y^4}{16y^5}
       =\frac7{8y}+\frac5{16y^3}+\frac1{16y^5}.
\end{equation}
Its raw remainder scale is $(y+\operatorname{Re}y)^{-4}$ and its domain is the
same positive-core-coordinate condition as above. On the positive real ray,
\begin{equation}\label{eq:smallscale}
 R_2^-(y)=\frac1{16y^4}.
\end{equation}
The true small root is
\[
 g_-(y)=\frac1y+\frac1{y^3}+\frac2{y^5}+\frac5{y^7}+\cdots,
\]
so
\begin{equation}\label{eq:smallerror}
 g_-(y)-B_2(y)=\frac1{8y}+\frac{11}{16y^3}
                    +\frac{31}{16y^5}+\OO{y^{-7}}.
\end{equation}
Therefore
\begin{equation}\label{eq:smallratio}
 \frac{|g_-(y)-B_2(y)|}{R_2^-(y)}\sim 2y^3.
\end{equation}
The finite approximation misses one eighth of the leading inverse coefficient,
yet the object claims fourth-power decay of its error. Both this and the
large-source example are recorded under N01 because they share the same
admission and majorant mechanism.

\subsection{Exact numerical enclosures}
Table~\ref{tab:samples} uses an independent rational oracle, not additional
package executions. For each rational $y$, integer square-root arithmetic
constructs outward rational endpoints for $\sqrt{y^2-4}$ on a grid with 100
decimal places. Equation~\eqref{eq:roots} then encloses the selected root and
the true approximation error. The complete fraction endpoints are retained in
\file{evidence/exact-enclosures.json}; the displayed decimals are rounded
summaries, not themselves directed-rounding certificates.

\begin{table}[htbp]
\centering\small
\begin{tabular}{@{}lrrrr@{}}
\toprule
Branch & $y$ & True absolute error & Advertised scale & Error / scale\\
\midrule
Large & 5 & $5.77122\cdot10^{-2}$ & $10^{-2}$ & $5.77122$\\
Large & 20 & $1.26100\cdot10^{-2}$ & $6.25\cdot10^{-4}$ & $20.1760$\\
Large & 100 & $2.50088\cdot10^{-3}$ & $2.5\cdot10^{-5}$ & $100.035$\\
Large & 1000 & $2.50001\cdot10^{-4}$ & $2.5\cdot10^{-7}$ & $1000.0035$\\
\addlinespace
Small & 5 & $3.11922\cdot10^{-2}$ & $10^{-4}$ & $311.922$\\
Small & 20 & $6.33655\cdot10^{-3}$ & $3.90625\cdot10^{-7}$ & $16221.6$\\
Small & 100 & $1.25069\cdot10^{-3}$ & $6.25\cdot10^{-10}$ & $2.00110\cdot10^6$\\
Small & 1000 & $1.25001\cdot10^{-4}$ & $6.25\cdot10^{-14}$ & $2.000011\cdot10^9$\\
\bottomrule
\end{tabular}
\caption{Depth-two moving-core witnesses with $c=y$. Independent exact root
and error enclosures; decimals shown to a convenient precision.}
\label{tab:samples}
\end{table}

For the large root, the positive Catalan tail implies a ratio strictly greater
than $y$ at the tabulated points. For the small root, the leading mismatch and
positive subsequent coefficients give the cubic divergence. The proof is the
asymptotic calculation, while the rational intervals provide independently
checkable finite-target evidence.

\section{Source-level cause and the failed theorem premise}
\subsection{Validation observes different data from the algorithm}
The relevant source fragment is:
\begin{lstlisting}
validateInput[core + perturbation, limit];
If[x === y || ! FreeQ[core + perturbation, y],
  fail["InvalidVariables",
    "Use distinct source and target symbols, with no target symbol in the forward data."]];
...
model = corePerturbationModel[
  core, perturbation, x, coord, ell, ass];
\end{lstlisting}
The first two checks see the evaluated sum. The model sees the original
separate argument values. With
\[
 F_0=x-c,\quad H=c+x^{-1},\quad c=\operatorname{Re}y,
\]
the sum is target-free, but neither of the two consumed components is.
The corresponding finite-source split is
\[
 F_0=x^{-1}-c,\qquad H=c+x.
\]
No option grammar, native-backend fallback, series composition, supplied core
inverse, or numerical root finder is required to reach the defect.

\code{FreeQ} tests an expression for matching subexpressions
\cite{freeq}. It cannot recover a target occurrence already removed from its
argument by evaluation. The logical implication being assumed is invalid:
\begin{equation}\label{eq:invalidimplication}
 \operatorname{FreeQ}(F_0+H,y)
 \quad\not\Rightarrow\quad
 \operatorname{FreeQ}(F_0,y)\land\operatorname{FreeQ}(H,y).
\end{equation}
The tuple form preserves the distinction between the two argument values and
therefore tests the property actually needed by the algorithm.

\subsection{What the model computes}
Let $u$ be the positive source coordinate. In the large-source case $u=1/x$;
in the finite-source case $u=x$. In both cases, for a nonzero frozen $c$,
\[
 F_0(u)=u^{-1}-c,\qquad H(u)=c+u.
\]
The leading source power of the core is $p=-1$. The perturbation has powers
$0$ and $1$, so its minimum gap above the core is $\delta=1$. The core offset
is $-c$, and its local inverse coordinate is
\begin{equation}\label{eq:corecoordinate}
 u_0=\frac1{y+c}.
\end{equation}
For the default observable, the local observable power is $r_{\rm loc}=-1$
at the infinite source and $r_{\rm loc}=1$ at the finite source.

The implementation forms the truncation pair
\begin{equation}\label{eq:sourcebound}
 \bigl(r_{\rm loc}+(N+1)\delta,\;(N+1)d\bigr),
\end{equation}
where $d$ is a logarithmic-degree bound. Here $d=0$, so the stored scale is
$u_0^N$ in the large-source chart and $u_0^{N+2}$ in the small-source chart.
The source explicitly describes its majorant as a fixed-data existence result
based on fixed core and finite higher-power perturbation parameters
\cite{core}. The successful result does not provide uniform estimates for
$c=c(y)$.

\subsection{Why a correct marker calculation is not enough}
The marker algorithm differentiates with respect to the source coordinate,
holding coefficient parameters fixed. Its finite coefficients are valid for
that frozen-parameter formal calculation. After the target-dependent split is
admitted, however, those coefficients are specialized along a moving path in
parameter space. In these witnesses the supposedly constant term $c$ is of
order $y$, not of fixed magnitude.

Equivalently, the normalized constant perturbation is comparable to
$c/(y+c)$. When $c=y$, this is $1/2$, not a quantity tending to zero. Each marker
order can reduce a coefficient geometrically without improving its power of
$y$. An implicit-function theorem at fixed parameters does not imply the
power remainder claimed after that substitution. Section~\ref{sec:alldepth}
makes this precise without relying on numerical experimentation.

\subsection{Severity and containment}
N01 is a high-priority mathematical correctness issue because a successful
analytic result carries a false error class. Its trigger is an unsupported
moving split rather than an ordinary fixed decomposition. That makes the
repair small, but it does not make the current success sound: the constructor
must enforce the premise it advertises.

The report does not assert that a particular downstream interval certificate
has accepted the bad result or that every runtime reaches the same branch.
The demonstrated failure is already the returned remainder itself. The native
witness and exact inverse proof suffice to establish that narrower conclusion.

\section{Exact all-depth analysis}\label{sec:alldepth}
\subsection{Freeze the parameter before introducing the moving path}
Initially let $y>2$ and $c>0$ be separate fixed real numbers, and write
$A=y+c$. Introduce a marker $\eps$ multiplying the perturbation. The large-source
split gives
\[
 X-c+\eps(c+X^{-1})=y,
\]
or
\begin{equation}\label{eq:marklarge}
 X^2-(A-\eps c)X+\eps=0,\qquad X(0)=A.
\end{equation}
The finite-source split gives
\[
 U^{-1}-c+\eps(c+U)=y,
\]
or
\begin{equation}\label{eq:marksmall}
 \eps U^2-(A-\eps c)U+1=0,\qquad U(0)=A^{-1}.
\end{equation}
These equations independently encode the frozen-parameter marker problem
that the source recurrence solves. They need no implementation of
\code{GeneralizedSeries} to analyze.

Their analytic branches at $\eps=0$ are
\begin{align}
 X(\eps)&=\frac{A-\eps c+\sqrt{(A-\eps c)^2-4\eps}}2,\label{eq:exactmarkerlarge}\\
 U(\eps)&=\frac{2}{A-\eps c+\sqrt{(A-\eps c)^2-4\eps}}.\label{eq:exactmarkersmall}
\end{align}
In particular,
\begin{equation}\label{eq:reciprocalmarker}
 X(\eps)U(\eps)=1.
\end{equation}
At $\eps=1$, the parameter $c$ cancels from the equation, leaving precisely
$g_+$ and $g_-$. The finite Taylor polynomials, however, still depend on the
chosen decomposition.

\subsection{General coefficient formulas}
Define the Catalan numbers by
\[
 C_m=\frac1{m+1}\binom{2m}{m},\qquad m\geq0.
\]
The generating identity
\[
 \frac{1-\sqrt{1-4z}}{2z}=\sum_{m\geq0}C_mz^m
\]
follows by expanding the square root or solving its quadratic formal equation.
The following formulas are exact identities for Taylor coefficients; they are
not numerical fits.

\begin{proposition}[Frozen-parameter marker coefficients]\label{prop:coeffs}
Write $X(\eps)=\sum_{k\geq0}L_k\eps^k$ and
$U(\eps)=\sum_{k\geq0}S_k\eps^k$, with the branches above. Then $L_0=A$, and
for $k\geq1$,
\begin{equation}\label{eq:Lk}
 L_k=-c\,\mathbf1_{k=1}
 -\sum_{m=1}^{k}C_{m-1}\binom{k+m-2}{k-m}
                    \frac{c^{k-m}}{A^{k+m-1}}.
\end{equation}
For every $k\geq0$,
\begin{equation}\label{eq:Sk}
 S_k=\sum_{m=0}^{k}C_m\binom{k+m}{k-m}
                   \frac{c^{k-m}}{A^{k+m+1}}.
\end{equation}
\end{proposition}
\begin{proof}
Put $B=A-\eps c$. Applying the Catalan identity to
\eqref{eq:exactmarkerlarge}--\eqref{eq:exactmarkersmall} gives the formal series
\[
 X(\eps)=B-\sum_{m\geq1}\frac{C_{m-1}\eps^m}{B^{2m-1}},
 \qquad
 U(\eps)=\sum_{m\geq0}\frac{C_m\eps^m}{B^{2m+1}}.
\]
For a positive integer $j$,
\[
 (A-\eps c)^{-j}
 =\sum_{\ell\geq0}\binom{j+\ell-1}{\ell}
       \frac{c^\ell}{A^{j+\ell}}\eps^\ell.
\]
Taking the coefficient of $\eps^k$ in each expression gives
\eqref{eq:Lk} and \eqref{eq:Sk}. The extra $-c$ occurs only in $L_1$ because
$B=A-\eps c$ is linear in the marker.
\end{proof}

The initial coefficients illustrate the calculation:
\begin{align*}
 L_0&=A,& L_1&=-c-A^{-1},& L_2&=-cA^{-2}-A^{-3},\\
 S_0&=A^{-1},& S_1&=cA^{-2}+A^{-3},&
 S_2&=c^2A^{-3}+3cA^{-4}+2A^{-5}.
\end{align*}
Their sums through degree two, with $c=y$, give exactly the two native finite
expressions \eqref{eq:largeobserved} and \eqref{eq:smallobserved}.
The independent Python implementation verifies the quadratic equations and
reciprocal coefficient convolution at several exact parameter choices. The
proof above, not those finitely many tests, establishes the general formulas.

\subsection{Specialize to a moving core}
Only now set
\begin{equation}\label{eq:theta}
 c=\lambda y,\qquad \lambda>0,\qquad
 \theta=\frac{\lambda}{1+\lambda}\in(0,1).
\end{equation}
Define the depth-$N$ approximations at marker value one by
\[
 A_N(y)=\sum_{k=0}^{N}L_k\big|_{c=\lambda y},\qquad
 B_N(y)=\sum_{k=0}^{N}S_k\big|_{c=\lambda y}.
\]
The following is an all-depth deduction from the inspected recurrence and
independent quadratic equations. The public native observations themselves
were the $\lambda=1$, $N=2$ cases.

\begin{theorem}[Wrong orders at every fixed affected depth]\label{thm:orders}
Fix $\lambda>0$. For each fixed $N\geq1$,
\begin{equation}\label{eq:largeall}
 A_N(y)-g_+(y)=\frac{\theta^N}{y}+\OO{y^{-3}}.
\end{equation}
For each fixed $N\geq0$,
\begin{equation}\label{eq:smallall}
 g_-(y)-B_N(y)=\frac{\theta^{N+1}}{y}+\OO{y^{-3}}.
\end{equation}
The scales inferred by the current fixed-gap construction for nonzero $c$ are
$A^{-N}$ and $A^{-(N+2)}$, respectively. Their error ratios satisfy
\begin{align}
 \frac{|A_N-g_+|}{A^{-N}}&\sim\lambda^N y^{N-1},\qquad N\geq1,\label{eq:largeratioall}\\
 \frac{|g_--B_N|}{A^{-(N+2)}}&\sim
       \lambda^{N+1}(1+\lambda)y^{N+1},\qquad N\geq0.\label{eq:smallratioall}
\end{align}
Thus the large-source scale is false for every $N\geq2$, and the finite-source
scale is false for every $N\geq0$.
\end{theorem}
\begin{proof}
For fixed $k$, the $m=1$ term in \eqref{eq:Lk} is
\[
 -\frac{c^{k-1}}{A^k}
 =-\frac1{(1+\lambda)y}\theta^{k-1}.
\]
Every term with $m\geq2$ is $O(y^{-3})$, since its power of $y$ is
$1-2m$. For $N\geq1$, the $A-c$ contribution is $y$. Summing the geometric
progression yields
\[
 A_N(y)=y-\frac{1-\theta^N}{y}+O(y^{-3}).
\]
Subtract \eqref{eq:catalanroot} to obtain \eqref{eq:largeall}.

In \eqref{eq:Sk}, the $m=0$ term is
$c^k/A^{k+1}=\theta^k/((1+\lambda)y)$. All $m\geq1$ terms are $O(y^{-3})$.
Therefore
\[
 B_N(y)=\frac{1-\theta^{N+1}}{y}+O(y^{-3}),
\]
which proves \eqref{eq:smallall}. Finally $A=(1+\lambda)y$ and the source gap
is one, so division by the two advertised scales gives
\eqref{eq:largeratioall}--\eqref{eq:smallratioall}.
\end{proof}

\begin{remark}[Depth and zero-parameter exceptions]
For the large-source branch at $N=1$, the ratio tends to $\lambda$, so this
particular scale is not refuted at that depth. At $N=0$, the large approximation
is $(1+\lambda)y$ and its error is $\lambda y+O(y^{-1})$, while the inferred
scale is one; that depth also fails. These cases should not be blurred into an
incorrect claim about all depths having the same behavior.

When $c=0$, the constant perturbation row disappears. The minimum source gap is
then two, not one. The valid unshifted construction consequently uses scales
$y^{-(2N+1)}$ for the large branch and $y^{-(2N+3)}$ for the small branch.
For example, the patched native large-source depth-two control returns
$y-y^{-1}-y^{-3}$ with scale $y^{-5}$. Substituting $\lambda=0$ into a
nonzero-row analysis would miss this structural change.
\end{remark}

\subsection{Convergence in the marker does not repair the power claim}
There is a particularly instructive distinction here. For fixed $c>0$ and
$A>0$, the nearest positive discriminant zero of the exact marker functions is
\begin{equation}\label{eq:radius}
 \eps_- = \frac{\bigl(\sqrt{Ac+1}-1\bigr)^2}{c^2}.
\end{equation}
This follows by solving $(A-\eps c)^2-4\eps=0$; both branch points are positive,
and the smaller one is \eqref{eq:radius}. The branch at zero is analytic up to
that nearest square-root singularity. Under $c=\lambda y$,
\[
 \eps_-\longrightarrow\frac{1+\lambda}{\lambda}>1
 \quad(y\to\infty).
\]
For the actual $c=y$ witnesses, $\eps_->1$ whenever $y>2$. Hence their Taylor
series in the marker genuinely converge at $\eps=1$ on the ray used in the
counterexample.

There is no contradiction with Theorem~\ref{thm:orders}. Sending the marker
depth to infinity at a fixed target and sending the target to infinity at a
fixed marker depth answer different questions. A residual leading coefficient
of size $\theta^N$ can shrink geometrically with $N$ but remain attached to
$y^{-1}$. This example therefore isolates an invalid asymptotic grading, not
necessarily a divergent marker algorithm.

One can also see the warning in the first omitted large-source marker
coefficient. Formula~\eqref{eq:Lk} gives, at $c=y$,
\begin{equation}\label{eq:firstomitted}
 L_3=-\frac1{8y}-\frac3{16y^3}-\frac1{16y^5}.
\end{equation}
That term alone is not $O(y^{-2})$. Equation~\eqref{eq:firstomitted} is a
source-derived independent formula; a later native attempt to project the
stored omitted-term field failed at the service level, so no native observation
of that field is claimed.

\section{Candidate repair and its precise contract}
\subsection{The two-line change}
The supplied canonical-source patch is:
\begin{lstlisting}
- validateInput[core + perturbation, limit];
- If[x === y || ! FreeQ[core + perturbation, y],
+ validateInput[{core, perturbation}, limit];
+ If[x === y || ! FreeQ[{core, perturbation}, y],
    fail["InvalidVariables", ...]];
\end{lstlisting}
The second changed line is the essential soundness repair for N01. With an
ordinary list preserving both values, \code{FreeQ} can see a target occurrence
in either component even if the sum would cancel it. This uses the existing
structural policy; it does not attempt a new lexical free-variable analyzer
for held functions or binders.

The first changed line checks exactness, nonfinite content, and the generic
input budget on the same component tuple that the algorithm will consume.
It is a consistency hardening alongside the target guard, not a second
independently demonstrated defect. In particular, the native baseline already
rejects the attempted inexact cancellation control. That observation must not
be recast as proof of a new inexact-admission hole.

Tuple validation can charge a slightly different leaf count from sum
validation, especially where large subexpressions cancel. This is appropriate
for an algorithm that retains the separate components, but it is a budget-policy
change worth acknowledging. The minimal target-only variant leaves the first
line unchanged and still closes the reproduced counterexamples. The tested
candidate and supplied diff use both replacements.

\subsection{Preserve intentional dependencies}
The target is not forbidden from every argument. A supplied
\code{"CoreInverse" -> phi[y]} is expected to be a function of the target;
the existing source-independence and branch-identity checks for that expression
remain separate. A repair that adds \code{requested} to the target-free tuple
would reject legitimate use.

The patched native control
\begin{lstlisting}
AsymptoticAnalysis`AsymptoticCoreInverse[
  x - 1, 1 + 1/x, {x, Infinity}, {y, 2},
  "CoreInverse" -> y + 1
]
\end{lstlisting}
returns a successful analytic result with finite expression
\[
 y-\frac1{y+1}-\frac1{(y+1)^2}-\frac1{(y+1)^3}
\]
and remainder scale $(y+1)^{-2}$. The native difference from this expected
expression simplified exactly to zero. Here the core offset is truly fixed.
This is not the moving $c=y$ family.

\subsection{Patch-utility behavior}
\file{code/patch_core_validation.py} emits a unified diff by default. It checks
that each old literal anchor occurs exactly once, and refuses a missing or
duplicate anchor or an already-patched file. Its optional \code{--output}
argument creates a separate, previously nonexistent file; it never overwrites
the input or an existing destination. Eight fixture tests exercise these
properties. They are tests of the utility, not tests of the package.

The utility does not prove the input file's Git provenance, install the
package, edit the user's checkout, or rebuild the standalone distribution.
Use the pinned revision and inspect the emitted diff. After integrating the
canonical module change, regenerate the root standalone with the repository's
maintained builder \cite{sourcemap}. Maintaining a hand-edited generated fork
would make later source reconciliation unnecessarily difficult.

\subsection{Native patch observations}
The execution service imported the pinned standalone as text. Both replacement
anchors occurred once. Two literal replacements were applied, the result was
written to a temporary file, and a fresh native kernel loaded that file.
No canonical repository source was changed. Across two successful queries,
the following five public requests were observed:

\begin{center}
\begin{tabularx}{\textwidth}{@{}p{0.10\textwidth}p{0.43\textwidth}Y@{}}
\toprule
Case & Request class & Observed patched outcome\\
\midrule
CV01 & Large-source cancelling \code{Re[y]} &
\code{Failure["InvalidVariables",\ldots]}\\
CV03 & Finite-source cancelling \code{Re[y]} &
\code{Failure["InvalidVariables",\ldots]}\\
CV08 & Unshifted large source: core $x$, perturbation $1/x$, depth two &
$y-y^{-1}-y^{-3}$, scale $y^{-5}$\\
CV10 & Fixed offset with explicit \code{CoreInverse -> y+1} &
Expected expression preserved; scale $(y+1)^{-2}$\\
CV12 & Inexact-component negative control, depth one &
\code{Failure["InexactInput",\ldots]}\\
\bottomrule
\end{tabularx}
\end{center}
These are five focused observations, not an executed twelve-test MUnit file,
a full-suite result, or Mathics acceptance. Their manual transcription is in
\file{evidence/native-observations.json}.

\section{Validation record and limits}
\subsection{Native environment and failures}
The successful native service reports
\begin{lstlisting}
15.0.0 for Linux x86 (64-bit) (May 6, 2026)
\end{lstlisting}
The service is stateless across queries; each package query loaded its own
commit-pinned entry. Several connection and HTTP-502 attempts returned no usable
result. Those failures are excluded from every observation count. They neither
prove that a package call failed nor establish that its expected output occurred.

The durable native record is a manual transcription of successful connector
responses, not a machine-signed or automatically replayed receipt. It preserves
that provenance explicitly. No successful Mathics session or local Wolfram
executable was available, and the full repository test suite was not run.

\subsection{A separate native forward screen}
Seventeen ordinary forward inputs were expanded with
\begin{lstlisting}
s = AsymptoticAnalysis`AsymptoticExpansion[
  f, {x, 0, 4}, "Backend" -> "Package"
];
FullSimplify[Normal[Series[f - Normal[s], {x, 0, 3}]]]
\end{lstlisting}
Every input returned a \code{GeneralizedSeries}, and all seventeen displayed
differences were zero. The inputs cover pole-amplified Taylor data,
cancellation before a singular quotient, bounded composite arguments, exact
irrational powers, and varying powers such as $(1+x)^{1/x}$. Appendix~\ref{app:screen}
lists them with their returned remainder powers.

This is useful positive evidence for the exercised coefficient paths. It does
not certify every higher remainder order the objects reported: the comparison
was through powers below four. It is also not an independent implementation
of the native algebra system. It compares the package's output against the
native series of its difference with the source.

One additional observable control was successful:
\begin{lstlisting}
s = AsymptoticAnalysis`AsymptoticExpansion[
  Sin[x], {x, 0, 4}, "Backend" -> "Package"];
AsymptoticAnalysis`SeriesObservable[s, z/x^4, z]
(* finite expression: 1/x^3 - 1/(6 x)
   remainder: PowerLogRemainder[x, 1, 0] *)
\end{lstlisting}
Thus the varying multiplier transports the sine remainder to the expected
first power in this case. A suspicion that it was always treated as a constant
was not retained.

\subsection{Independent exact tests}
The Python companion uses only the standard library. It does not call Wolfram,
Mathics, SymPy, or the package. Its run completed with \textbf{30 test methods
passed, zero failures, zero errors, zero skips} on Python 3.13.5. These split
into 22 exact mathematical checks and eight patch-utility fixture checks.

The mathematical checks include the Catalan coefficient formulas, the two
implicit quadratic equations through order ten at selected exact parameters,
and the reciprocal convolution \eqref{eq:reciprocalmarker} through order
eight. They compare the depth-two formulas with the native observations, test
the unshifted Catalan specialization and its changed source gap, verify
outward root brackets, and examine the predicted leading-error constants at
several fixed depths and positive rational $\lambda$. The all-depth tests at a
large exact target are finite checks of a proved formula, not proofs of a
limit by sampling.

To enclose a rational square root $\sqrt{a/b}$, the oracle sets $D=10^d$ and
\[
 k=\operatorname{isqrt}\!\left(\left\lfloor aD^2/b\right\rfloor\right).
\]
Then $k/D\leq\sqrt{a/b}<(k+1)/D$, except that an exact equality is retained when
$(k/D)^2=a/b$. These inequalities use integers and rational arithmetic only.
The small-root formula $2/(y+\sqrt{y^2-4})$ avoids subtracting close endpoints.
The resulting rational error enclosures are independent of floating-point
roundoff and of the package's own remainder claims.

The oracle is deliberately narrow. Its bounded depth and rational-input
checks are engineering safeguards, not a general resource sandbox or a new
asymptotic package. It is an executable mathematical counterexample and
regression reference for this audit.

\subsection{Supplied acceptance files versus executed observations}
The archive includes a twelve-case desired post-patch MUnit file and a
four-case no-MUnit portable probe. Neither complete file was executed during
this audit. Some corresponding requests were executed inline, as specified in
the patch table, but that does not make the whole file an acceptance receipt.

The twelve cases add controls for \code{Abs[y]}, plain \code{y}, a target in
only one component, identical source/target variables, a regular finite-endpoint
example, and a fixed real symbolic parameter. The four portable cases focus on
the two bad splits, the unshifted control, and the legitimate supplied inverse.
Each file requires the package to have been loaded first. Public package names
are fully qualified, and the portable file avoids dependence on MUnit.

\begin{center}
\begin{tabularx}{\textwidth}{@{}p{0.43\textwidth}Y@{}}
\toprule
Evidence category & What was actually established\\
\midrule
Public baseline counterexamples & Two native successful objects; exact analysis
refutes their returned remainder scales.\\
Patched native requests & Five requested outcomes observed in temporary patched
standalone copies.\\
Native forward screen & Seventeen coefficient differences below power four
vanish; no exhaustive remainder certification.\\
Independent Python tests & Thirty methods pass; none loads the package.\\
Supplied MUnit / portable files & Twelve / four desired cases supplied;
complete files not executed.\\
Mathics / full repository suite & Not executed.\\
\bottomrule
\end{tabularx}
\end{center}
No pooled success total is meaningful across these categories.

\section{Wolfram 15 and Mathics3 implications}
\subsection{What transfers at the source level}
The defective guard is in shared canonical source, not solely in a Wolfram
backend adapter. Both public argument values must be target-free for the
fixed-data exact-core construction, regardless of which evaluator eventually
executes it. The proposed list-based check uses the same existing
\code{FreeQ} policy and ordinary list structure rather than a runtime-specific
symbolic solver. This makes it a natural shared-source repair.

It does not follow that the unpatched Mathics package must return the same bad
objects. Earlier realness, sign, automatic-core, or simplification behavior
may refuse a candidate before the faulty majorant is exposed. Such a refusal
would be safe for that request, not proof that the shared missing guard is
correct. The correct Mathics statement from this audit is therefore
\emph{shared source obligation, fresh public execution unverified}.

\subsection{A focused two-runtime acceptance target}
For this repair, the decisive behavior is simple: both target-cancelling splits
must be rejected by the dependency check, and fixed-data controls must remain
successful where that runtime supports the underlying operations. The supplied
portable probe provides four concrete cases rather than treating package load
as a compatibility claim. The native MUnit file adds a wider twelve-case
contract matrix.

The modular and generated standalone entries should both contain the repaired
guard after rebuilding. Run the same focused cases in fresh Wolfram 15 and
Mathics sessions, record the actual runtime/version and entry, and retain a
structured distinction between the intended dependency refusal and an earlier
unrelated capability refusal. This is an acceptance obligation for N01, not a
reissued recommendation for a wholesale new CI or compatibility architecture.

The user guide already describes complete Mathics support and complete native
expansion coverage as unfinished goals \cite{guide}. This report does not add
those familiar gaps as new findings, and it does not infer Windows 15.0.1
behavior from the Linux 15.0.0 observations.

\subsection{Performance and resource interpretation}
No new end-to-end timing or memory defect is claimed. The target-free tuple
check traverses the actual retained inputs rather than relying on their
possibly simplified sum. Its intended benefit is sound admission, not a
measured speedup. It can also avoid spending downstream symbolic work on an
invalid split, but no benchmark quantifies that effect here.

The independent reference implementation is fast for these small finite
quadratics; that timing says nothing about package performance. Existing
reports already cover dense sparse-polynomial allocations, repeated symbolic
work, recurrence termination, and runner resource concerns. Their
recommendations are excluded rather than relabeled as opportunities found in
this audit \cite{register,wave4,wave6}.

\section{Specific API and development follow-through}
\subsection{Make the rejected component explicit}
The present failure message says that forward data must not contain the
target, but users can see a target-free stored sum and reasonably believe that
condition holds. After the repair, a useful focused diagnostic improvement is
to identify which positional component contains the forbidden target.
For example, retain \code{Failure["InvalidVariables",\ldots]} while adding
\code{"TargetVariable" -> y} and
\code{"DependentComponents" -> \{"Core", "Perturbation"\}}.
A compact component-name list is preferable to duplicating a potentially large
expression in the failure data.

This diagnostic is a proposal, not part of the two-line patch or its five
native observations. It should be derived from the same component values used
by the guard, without an extra simplifying sum or repeated evaluation of user
expressions. The current tag can remain stable for callers that dispatch on
it.

\subsection{Document three different dependency roles}
The exact-core section of the guide should state the newly enforced boundary
in concrete terms: each of the core and perturbation arguments must be
independent of the requested target; cancellation between them does not satisfy
that condition. A supplied inverse is instead expected to depend on the target
but not on the eliminated source. Parameter assumptions refer to fixed
parameters under the existing policy.

A useful documentation pair is the refused moving split
\code{x-Re[y], Re[y]+1/x} beside the accepted fixed split
\code{x-1, 1+1/x} with \code{"CoreInverse"->y+1}. That example explains the
actual distinction much better than a blanket statement that all auxiliary data
must be target-free. This is a proposed clarification specific to N01, not a
claim that every related API should receive identical restrictions.

\subsection{Use sum-preserving input transformations as tests}
The mechanism suggests a small, reproducible regression family. Starting with
a valid decomposition $(F_0,H)$, replace it by
\begin{equation}\label{eq:metamorphic}
 (F_0-C,\;H+C).
\end{equation}
The forward sum is unchanged, but the marker decomposition is changed. With a
fixed admissible constant $C$, the constructor may legitimately return a
different finite marker approximation with its corresponding bound. Equality
of finite approximations at equal marker depth must \emph{not} be imposed.
With $C=\operatorname{Re}y$ or another retained target-dependent expression,
the current API should reject the transformed request before theorem-based
construction.

This tests the information-loss boundary directly. It distinguishes a true
sum-invariant property---the underlying forward equation---from a
split-dependent property---finite marker coefficients and their proof
premises. The supplied native tests instantiate both sides of that distinction.

\subsection{Retain the reciprocal marker oracle}
For the two quadratic charts, the coefficient identity
\[
 \sum_{j=0}^k L_jS_{k-j}=\begin{cases}1,&k=0,\\0,&k>0\end{cases}
\]
is an exact cross-chart check. Together with each defining quadratic equation,
it can detect a recurrence or observable reconstruction error independently of
public formatting and remainder metadata. It is already implemented in the
companion exact tests.

This invariant is deliberately restricted to the stated paired family. It is
not proposed as a universal inverse-residual checker or an excuse to accept
moving coefficients. Passing the algebraic identity proves a property of the
marker coefficients, while target-independence remains a separate admission
obligation.

\subsection{Moving cores are outside this repair}
The exact formulas show that a moving core can be algebraically meaningful and
can even have a convergent marker series. Supporting it with useful asymptotic
error statements would nevertheless require an explicit joint grading or a
uniform parameter-dependent bound. A realness proof, successful simplification,
or small numerical marker residual is not that theorem.

For the current API, reject the unsupported moving split. Replacing the result
by an automatically chosen different core would change the user's requested
marker decomposition and should not happen silently. A future explicitly
separate mode could choose a different contract, but this report makes no claim
that such a mode has been designed, implemented, or validated. The immediate
repair should remain small and auditable.

\section{Reproduction and archive guide}
\subsection{Rebuild the article}
The article is self-contained LaTeX using standard text, mathematics, table,
and listing packages. No external figure or font file is required:
\begin{lstlisting}
cd article
pdflatex -interaction=nonstopmode -halt-on-error review.tex
pdflatex -interaction=nonstopmode -halt-on-error review.tex
# Repeat once more if cross-reference warnings remain.
\end{lstlisting}
The delivered PDF was compiled from this source. The build summary and layout
checks are recorded separately in the evidence directory.

\subsection{Run the independent mathematics and patch tests}
From the extracted archive root:
\begin{lstlisting}
python code/run_checks.py
\end{lstlisting}
Python 3.10 or later is required by the companion type syntax; the recorded run
used 3.13.5. There are no third-party Python dependencies. The command writes
its transcript, summary, exact enclosures, and decimal sample table into
\file{evidence/}. It does not execute or download the Wolfram package. Those
new local outputs replace the companion run records, not repository files.

\subsection{Prepare the canonical patch}
With a checkout of the pinned source:
\begin{lstlisting}
python code/patch_core_validation.py \
  /path/to/Asymptotic/src/Kernel/CorePerturbation.wl
\end{lstlisting}
The default output is a diff for inspection. To create a distinct candidate
file without overwriting the source:
\begin{lstlisting}
python code/patch_core_validation.py \
  /path/to/Asymptotic/src/Kernel/CorePerturbation.wl \
  --output /new/path/CorePerturbation.patched.wl
\end{lstlisting}
The output path must not already exist. Integrate the reviewed changes into the
canonical module using the usual repository workflow, then run the maintained
standalone builder from the repository root:
\begin{lstlisting}
python validation/build_standalone.py
\end{lstlisting}
No such integration into the user's repository was performed in this audit.

\subsection{Run package observations and acceptance}
In a fresh Wolfram kernel, first load a selected local package entry, then
load the observation file or run the post-patch MUnit file as separate inputs:
\begin{lstlisting}
Get["/path/to/AsymptoticAnalysis.wl"];
Get["/path/to/audit/code/ProbeComponents.wl"];

(* For the patched package: *)
TestReport["/path/to/audit/code/ReviewComponents.wlt"]
\end{lstlisting}
For a Mathics session, load the package first, using its maintained runtime
instructions and evaluator budget, then evaluate:
\begin{lstlisting}
Get["/path/to/audit/code/PortableComponents.wl"];
\end{lstlisting}
Do not treat this instruction as a recorded successful Mathics run. The probe
prints a runtime-tagged result with passed and failed case counts. It contains
no automatic network download, repository edit, or global replacement of
package definitions.

\subsection{Files and provenance}
\begin{longtable}{@{}p{0.40\textwidth}p{0.54\textwidth}@{}}
\toprule
File or group & Purpose\\
\midrule\endhead
\file{article/review.tex}, \file{review.pdf} & Complete article source and
compiled report.\\
\file{code/component-validation.diff} & Two-line canonical-source candidate.\\
\file{code/patch_core_validation.py} & Checked, non-overwriting diff/copy utility.\\
\file{code/reference.py} & Independent exact quadratic/marker/root oracle.\\
\file{code/test_independent.py}, \file{run_checks.py} & Thirty independent
mathematics and patch-fixture test methods and evidence generator.\\
\file{code/ReviewComponents.wlt} & Twelve desired post-patch native acceptance
cases; complete file unexecuted in the audit.\\
\file{code/PortableComponents.wl} & Four no-MUnit regression cases; Mathics
execution unverified.\\
\file{code/ProbeComponents.wl} & Public observation probe for baseline or
patched source.\\
\file{code/ForwardScreen.wl} & Reproducible case list for the separately
executed native coefficient screen.\\
\file{evidence/native-observations.json} & Manual transcription of successful
native responses and explicit execution limits.\\
\file{evidence/exact-enclosures.json}, \file{numerical-samples.csv} & Independent
rational root/error endpoints and rounded decimal summaries.\\
\file{evidence/independent-tests.*} & Actual Python test transcript and summary.\\
\file{evidence/novelty-ledger.csv}, \file{source-inventory.csv}, \file{scope.json}
& Non-overlap rationale, inspected-source coverage, and provenance boundaries.\\
\bottomrule
\end{longtable}
The ZIP contains no checksum files, full repository mirror, bundled third-party
libraries, font files, or TeX auxiliary build products. Commit identifiers
identify the source baseline; they are not standalone checksum files.

\section{Conclusion}
The new defect is an information-loss error at an admission boundary: validation
observes an evaluated sum, whereas the algorithm and its theorem depend on the
separate summands. The two public counterexamples are elementary enough to
exclude uncertainty about special-function conventions, numerical solvers,
branch selection, or hidden constants. Their exact inverses make the false
remainder statements unavoidable.

The corrective action is equally specific. Preserve the components while
checking their target independence, preserve the distinct dependency role of a
supplied inverse, regenerate the standalone from canonical source, and execute
the focused component tests on both supported runtimes. The tested candidate
already rejects both native witnesses and preserves selected valid controls;
full Mathics and broader acceptance remain unexecuted obligations, not assumed
successes.

Finally, the all-depth calculation explains why a plausible formal computation
can still produce an invalid analytic object. Convergent marker coefficients,
a target-free final equation, and a real local coordinate do not by themselves
supply a fixed-data asymptotic remainder. Enforcing that missing premise is the
substantive contribution of this incremental audit.

\appendix
\section{The native forward-screen inputs}\label{app:screen}
All cases used an exclusive package cutoff of four at $x\to0^+$ and an
explicit \code{"Backend"->"Package"}. All seventeen native comparisons of
the source minus the finite approximation through power three simplified to
zero. The final column is the \emph{returned} remainder power; this screen did
not independently prove every such stronger remainder. All returned logarithmic
degrees were zero.

\begin{longtable}{@{}rp{0.72\textwidth}r@{}}
\toprule
Case & Source expression & Power\\
\midrule\endhead
1 & $\sin x/x^2$ & 5\\
2 & $(e^x-1-x)/x^3$ & 4\\
3 & $\log(1+x)/x^4$ & 4\\
4 & $(\sqrt{1+x}-1-x/2)/x^3$ & 4\\
5 & $1/(\sin x/x-1)$ & 4\\
6 & $(1-\cos x)/(x-\sin x)$ & 5\\
7 & $\exp((\sin x-x)/x^3)$ & 4\\
8 & $\log((1-\cos x)/x^2)$ & 4\\
9 & $(1+\log(1+x)/x)^{-1}$ & 4\\
10 & $\sin(\log(1+x)/x)$ & 4\\
11 & $(1+\sin x/x)^{\sqrt2}$ & 4\\
12 & $((e^x-1)/x)^{1/x}$ & 4\\
13 & $(1+x)^{1/x}$ & 4\\
14 & $\exp(\sin x/x)$ & 4\\
15 & $\sqrt{1+x}-\sqrt{1-x}$ & 5\\
16 & $\log(1+x)+\log(1-x)$ & 4\\
17 & $\exp(\log(1+x))$ & $\infty$\\
\bottomrule
\end{longtable}

\section{Ruled-out and excluded candidates}
The guide explicitly defines the nonunit finite-endpoint observable power as
$(x-x_0)^r$, while power one reconstructs $x$ itself. The inspected Fourier
construction follows that convention. Treating it as a lost translation would
have been an incorrect new finding \cite{guide}.

The attempted inexact-component request was refused in the native baseline.
The tuple-validation patch therefore preserves and aligns exact-input checking;
it is not presented as fixing a newly reproduced approximate-number escape.
Likewise the ordinary observable multiplier example transported its remainder
correctly. Negative candidates and successful controls are part of the audit
record, not items to be silently discarded in order to increase the defect count.

Dense polynomial/Fourier support allocation, nonprincipal Mathics
\code{ProductLog}, previously described validation-runner issues, and the
nearby \code{SourceShift} error were already indexed. Their existing
recommendations were not repeated. This boundary is recorded both here and in
the novelty ledger so that the report remains an incremental contribution
rather than another copy of the established backlog.

\section{Acceptance case map}
\begin{longtable}{@{}lp{0.68\textwidth}@{}}
\toprule
ID & Desired post-patch contract\\
\midrule\endhead
CV01 & Refuse target dependence hidden by cancellation in the large-source split.\\
CV02 & Refuse the corresponding \code{Abs[y]} spelling.\\
CV03 & Refuse target dependence hidden by cancellation in the finite-source split.\\
CV04 & Diagnose plain target dependence before an unrelated coefficient-realness refusal.\\
CV05 & Refuse a target occurring only in the core.\\
CV06 & Refuse a target occurring only in the perturbation.\\
CV07 & Refuse identical source and target variables.\\
CV08 & Preserve the unshifted depth-two inverse and its fifth-power error scale.\\
CV09 & Preserve a regular finite-endpoint inverse marker expansion.\\
CV10 & Preserve a target-dependent supplied inverse for a fixed target-free core.\\
CV11 & Preserve a fixed real symbolic parameter under its assumptions.\\
CV12 & Preserve refusal of inexact input.\\
\bottomrule
\end{longtable}
Only CV01, CV03, CV08, CV10, and the depth-one CV12 request were observed inline
in patched native executions. The full twelve-case file and the four-case
portable file remain supplied acceptance artifacts, not recorded passing suites.

\begin{thebibliography}{99}
\small
\bibitem{repo} Vladimir Reshetnikov, \emph{Asymptotic repository}, pinned revision
\rev, 10 September 2026 (full identifier on the title page).
\href{https://github.com/VladimirReshetnikov/Asymptotic/tree/efa1aeec4845a9c35e140963a0333d0c9ec33b05}{[commit-pinned source]}.
\bibitem{core} \emph{CorePerturbation.wl}, same revision. Constructor admission at
lines 126--127; model, marker recurrence, and remainder contract in the same file.
\href{https://github.com/VladimirReshetnikov/Asymptotic/blob/efa1aeec4845a9c35e140963a0333d0c9ec33b05/src/Kernel/CorePerturbation.wl}{[commit-pinned source]}.
\bibitem{sourcemap} \emph{Kernel source map}, same revision.
\href{https://github.com/VladimirReshetnikov/Asymptotic/blob/efa1aeec4845a9c35e140963a0333d0c9ec33b05/src/Kernel/README.md}{[commit-pinned source]}.
\bibitem{guide} \emph{AsymptoticAnalysis User Guide}, same revision. Entry points,
Mathics instructions, power convention, and current coverage boundaries.
\href{https://github.com/VladimirReshetnikov/Asymptotic/blob/efa1aeec4845a9c35e140963a0333d0c9ec33b05/src/Documentation/UserGuide.md}{[commit-pinned source]}.
\bibitem{reviewindex} \emph{External code-review index}, same revision.
\href{https://github.com/VladimirReshetnikov/Asymptotic/blob/efa1aeec4845a9c35e140963a0333d0c9ec33b05/external-reports/code-review/README.md}{[commit-pinned source]}.
\bibitem{register} \emph{Code review implementation status}, same revision.
\href{https://github.com/VladimirReshetnikov/Asymptotic/blob/efa1aeec4845a9c35e140963a0333d0c9ec33b05/docs/development/CODE_REVIEW_STATUS.md}{[commit-pinned source]}.
\bibitem{wave3} \emph{Wave 3 intake}, same revision.
\href{https://github.com/VladimirReshetnikov/Asymptotic/blob/efa1aeec4845a9c35e140963a0333d0c9ec33b05/docs/development/WAVE_3_INTAKE.md}{[commit-pinned source]}.
\bibitem{wave4} \emph{Wave 4 intake}, same revision.
\href{https://github.com/VladimirReshetnikov/Asymptotic/blob/efa1aeec4845a9c35e140963a0333d0c9ec33b05/docs/development/WAVE_4_INTAKE.md}{[commit-pinned source]}.
\bibitem{wave5} \emph{Wave 5 review index}, same revision.
\href{https://github.com/VladimirReshetnikov/Asymptotic/blob/efa1aeec4845a9c35e140963a0333d0c9ec33b05/external-reports/code-review/wave-5/README.md}{[commit-pinned source]}.
\bibitem{wave6} \emph{Wave 6 review index}, same revision.
\href{https://github.com/VladimirReshetnikov/Asymptotic/blob/efa1aeec4845a9c35e140963a0333d0c9ec33b05/external-reports/code-review/wave-6/README.md}{[commit-pinned source]}.
\bibitem{r46} \emph{Code review 46: Target-dependent source shifts}, original report
at its own baseline, retained in the pinned tree. README and article consulted.
\href{https://github.com/VladimirReshetnikov/Asymptotic/blob/efa1aeec4845a9c35e140963a0333d0c9ec33b05/external-reports/code-review/wave-6/code-review-46/README.md}{[commit-pinned source]}.
\bibitem{r43} \emph{Code review 43: Focused source-review deltas}, original report
at its own baseline, retained in the pinned tree. F03 summarized in its README.
\href{https://github.com/VladimirReshetnikov/Asymptotic/blob/efa1aeec4845a9c35e140963a0333d0c9ec33b05/external-reports/code-review/wave-5/code-review-43/README.md}{[commit-pinned source]}.
\bibitem{freeq} Wolfram Research, \emph{FreeQ}, Wolfram Language Documentation,
accessed 10 September 2026.
\url{https://reference.wolfram.com/language/ref/FreeQ.html}.
\bibitem{re} Wolfram Research, \emph{Re}, Wolfram Language Documentation,
accessed 10 September 2026.
\url{https://reference.wolfram.com/language/ref/Re.html}.
\end{thebibliography}
\end{document}
